\documentclass[11pt]{article}

\usepackage[margin=0.9in]{geometry}
\usepackage{amsmath,amssymb}
\usepackage{booktabs}
\usepackage{enumitem}
\usepackage{listings}
\usepackage{xcolor}
\usepackage{hyperref}

\definecolor{backcolour}{rgb}{0.95,0.95,0.92}
\definecolor{codeblue}{rgb}{0.05,0.15,0.55}
\definecolor{codegreen}{rgb}{0.0,0.45,0.0}

\lstset{
    backgroundcolor=\color{backcolour},
    basicstyle=\ttfamily\small,
    keywordstyle=\color{codeblue}\bfseries,
    commentstyle=\color{codegreen},
    stringstyle=\color{purple},
    breaklines=true,
    frame=single,
    showstringspaces=false,
    tabsize=4
}

\title{CPSC 250 \\ NumPy, Plotting, and Simple Polynomial Fits}
\author{Edward Brash}
\date{}

\begin{document}

\maketitle

\section{The Big Idea}

So far, we have used Python lists as containers for data.

Lists are flexible and useful, but they are not designed primarily for numerical computation.

This week, we begin using tools that are designed for scientific and data-oriented computing.

\[
\text{Lists are containers.}
\]

\[
\text{NumPy arrays are mathematical objects.}
\]

This distinction is extremely important.

\section{Why NumPy?}

NumPy is a Python library for efficient numerical computing.

It allows us to store and manipulate large collections of numbers.

\begin{center}
\begin{tabular}{ll}
\toprule
Python List & NumPy Array \\
\midrule
General-purpose container & Numerical data structure \\
Stores arbitrary objects & Stores numbers efficiently \\
Often requires loops & Supports vectorized operations \\
Flexible & Fast \\
\bottomrule
\end{tabular}
\end{center}

\section{Importing NumPy}

The usual convention is:

\begin{lstlisting}[language=Python]
import numpy as np
\end{lstlisting}

This means that we access NumPy functions using:

\begin{lstlisting}[language=Python]
np.function_name()
\end{lstlisting}

\section{Creating a NumPy Array}

We can create a NumPy array from a Python list.

\begin{lstlisting}[language=Python]
import numpy as np

values = np.array([1, 2, 3, 4, 5])

print(values)
\end{lstlisting}

This prints:

\begin{verbatim}
[1 2 3 4 5]
\end{verbatim}

\section{Arrays Behave Differently from Lists}

Consider a Python list:

\begin{lstlisting}[language=Python]
numbers = [1, 2, 3]

print(numbers * 2)
\end{lstlisting}

This produces:

\begin{verbatim}
[1, 2, 3, 1, 2, 3]
\end{verbatim}

The list is repeated.

Now consider a NumPy array:

\begin{lstlisting}[language=Python]
numbers = np.array([1, 2, 3])

print(numbers * 2)
\end{lstlisting}

This produces:

\begin{verbatim}
[2 4 6]
\end{verbatim}

Each element is multiplied by 2.

\section{Special Arrays}

NumPy provides several useful functions for creating arrays.

\subsection{Zeros and Ones}

\begin{lstlisting}[language=Python]
zeros = np.zeros(5)
ones = np.ones(5)

print(zeros)
print(ones)
\end{lstlisting}

Output:

\begin{verbatim}
[0. 0. 0. 0. 0.]
[1. 1. 1. 1. 1.]
\end{verbatim}

\subsection{Using arange}

\begin{lstlisting}[language=Python]
x = np.arange(0, 10, 1)

print(x)
\end{lstlisting}

Output:

\begin{verbatim}
[0 1 2 3 4 5 6 7 8 9]
\end{verbatim}

The syntax is:

\[
\texttt{np.arange(start, stop, step)}
\]

\section{Using linspace}

For plotting and numerical work, \texttt{linspace} is often more useful.

\begin{lstlisting}[language=Python]
x = np.linspace(0, 10, 6)

print(x)
\end{lstlisting}

Output:

\begin{verbatim}
[ 0.  2.  4.  6.  8. 10.]
\end{verbatim}

The syntax is:

\[
\texttt{np.linspace(start, stop, number\_of\_points)}
\]

This creates evenly spaced values from the starting value to the ending value.

\section{Vectorized Operations}

Suppose we want to compute:

\[
y = x^2
\]

for many values of \(x\).

With a list, we might write a loop.

\begin{lstlisting}[language=Python]
x_values = [-2, -1, 0, 1, 2]
y_values = []

for x in x_values:
    y_values.append(x**2)

print(y_values)
\end{lstlisting}

With NumPy, we can write:

\begin{lstlisting}[language=Python]
x = np.array([-2, -1, 0, 1, 2])

y = x**2

print(y)
\end{lstlisting}

Output:

\begin{verbatim}
[4 1 0 1 4]
\end{verbatim}

No loop is required.

\section{Vectorized Physics Example}

Suppose an object moves according to:

\[
x(t) = x_0 + v_0t + \frac{1}{2}at^2
\]

We can calculate its position at many times.

\begin{lstlisting}[language=Python]
t = np.linspace(0, 10, 101)

x0 = 5.0
v0 = 12.0
a = -9.8

x = x0 + v0*t + 0.5*a*t**2
\end{lstlisting}

This calculation happens for every value in the array \texttt{t}.

This is the style of computation used in many scientific Python programs.

\section{Useful NumPy Functions}

NumPy also provides useful statistical and summary functions.

\begin{lstlisting}[language=Python]
data = np.array([72, 75, 79, 81, 78, 74])

print(np.min(data))
print(np.max(data))
print(np.mean(data))
print(np.std(data))
\end{lstlisting}

These compute the minimum, maximum, mean, and standard deviation.

\section{Plotting with Matplotlib}

We usually import Matplotlib as:

\begin{lstlisting}[language=Python]
import matplotlib.pyplot as plt
\end{lstlisting}

Suppose we have data:

\begin{lstlisting}[language=Python]
x = np.array([1, 2, 3, 4, 5])
y = np.array([3, 5, 7, 8, 11])
\end{lstlisting}

We can make a scatter plot:

\begin{lstlisting}[language=Python]
plt.scatter(x, y)

plt.xlabel("x")
plt.ylabel("y")

plt.show()
\end{lstlisting}

\section{Adding Labels and a Title}

Plots should be readable.

A plot should usually include:

\begin{itemize}
    \item an \(x\)-axis label,
    \item a \(y\)-axis label,
    \item a title,
    \item a legend, if more than one curve or data set appears.
\end{itemize}

\begin{lstlisting}[language=Python]
plt.scatter(x, y, label="Data")

plt.xlabel("x")
plt.ylabel("y")
plt.title("Example Data")

plt.legend()
plt.show()
\end{lstlisting}

\section{Linear Models}

Many data sets are approximately linear.

A straight line has the form:

\[
y = mx + b
\]

where:

\begin{itemize}
    \item \(m\) is the slope,
    \item \(b\) is the intercept.
\end{itemize}

In data analysis, we often want to find the line that best fits a set of data.

\section{Simple Fitting with polyfit}

NumPy includes a function called \texttt{polyfit}.

For a straight-line fit:

\begin{lstlisting}[language=Python]
m, b = np.polyfit(x, y, 1)
\end{lstlisting}

The final argument is the degree of the polynomial.

\[
1 \rightarrow \text{first-degree polynomial}
\]

A first-degree polynomial is a line:

\[
y = mx + b
\]

\section{Example: Finding the Best-Fit Line}

\begin{lstlisting}[language=Python]
import numpy as np
import matplotlib.pyplot as plt

x = np.array([1, 2, 3, 4, 5])
y = np.array([3, 5, 7, 8, 11])

m, b = np.polyfit(x, y, 1)

print("slope =", m)
print("intercept =", b)
\end{lstlisting}

Possible output:

\begin{verbatim}
slope = 1.9
intercept = 0.9
\end{verbatim}

So the best-fit line is approximately:

\[
y = 1.9x + 0.9
\]

\section{Plotting the Fit}

Once we have \(m\) and \(b\), we can compute the fitted values.

\begin{lstlisting}[language=Python]
y_fit = m*x + b
\end{lstlisting}

Then we plot both the data and the fit.

\begin{lstlisting}[language=Python]
plt.scatter(x, y, label="Data")

plt.plot(x, y_fit, label="Best-fit line")

plt.xlabel("x")
plt.ylabel("y")
plt.title("Data with Linear Fit")

plt.legend()
plt.show()
\end{lstlisting}

\section{A Smoother Fit Line}

The previous plot only draws the line at the original data points.

For a smoother-looking line, create more \(x\)-values.

\begin{lstlisting}[language=Python]
x_fit = np.linspace(min(x), max(x), 100)

y_fit = m*x_fit + b

plt.scatter(x, y, label="Data")
plt.plot(x_fit, y_fit, label="Best-fit line")

plt.xlabel("x")
plt.ylabel("y")
plt.legend()

plt.show()
\end{lstlisting}

\section{Example: Distance vs Time}

Suppose we measure the position of an object at several times.

\begin{center}
\begin{tabular}{cc}
\toprule
Time (s) & Distance (m) \\
\midrule
0 & 1.2 \\
1 & 3.1 \\
2 & 5.4 \\
3 & 7.0 \\
4 & 9.3 \\
5 & 11.1 \\
\bottomrule
\end{tabular}
\end{center}

We can store the data:

\begin{lstlisting}[language=Python]
time = np.array([0, 1, 2, 3, 4, 5])
distance = np.array([1.2, 3.1, 5.4, 7.0, 9.3, 11.1])
\end{lstlisting}

\section{Fitting the Distance Data}

\begin{lstlisting}[language=Python]
m, b = np.polyfit(time, distance, 1)

print("slope =", m)
print("intercept =", b)
\end{lstlisting}

In this example, the slope has units of:

\[
\frac{\text{meters}}{\text{second}}
\]

So the slope represents an estimate of the object's velocity.

\section{Plotting the Distance Data}

\begin{lstlisting}[language=Python]
time_fit = np.linspace(min(time), max(time), 100)

distance_fit = m*time_fit + b

plt.scatter(time, distance, label="Measurements")
plt.plot(time_fit, distance_fit, label="Best-fit line")

plt.xlabel("Time (s)")
plt.ylabel("Distance (m)")
plt.title("Distance vs Time")

plt.legend()
plt.show()
\end{lstlisting}

\section{Higher-Degree Polynomial Fits}

The same function can also fit higher-degree polynomials.

For example:

\begin{lstlisting}[language=Python]
coefficients = np.polyfit(x, y, 2)
\end{lstlisting}

This fits a quadratic function:

\[
y = ax^2 + bx + c
\]

However, today we will mostly focus on first-degree fits.

\section{A Warning About Polynomial Fits}

Just because we can fit a polynomial does not mean we should.

High-degree polynomials can wiggle in strange ways and may not represent the underlying situation.

For this course, our main goal is simple:

\[
\text{Use a fit to summarize a trend in the data.}
\]

\section{Practice}

Write a program that does the following:

\begin{enumerate}
    \item Creates NumPy arrays for time and distance.
    \item Makes a scatter plot of the data.
    \item Uses \texttt{np.polyfit()} to find the best-fit line.
    \item Prints the slope and intercept.
    \item Plots the data and the best-fit line together.
    \item Interprets the slope in words.
\end{enumerate}

\section{Summary}

Today we learned:

\begin{itemize}
    \item NumPy arrays are mathematical objects.
    \item Array operations are vectorized.
    \item \texttt{linspace} is useful for creating smooth ranges of values.
    \item Matplotlib lets us visualize numerical data.
    \item \texttt{np.polyfit()} can find a simple best-fit line.
    \item The slope and intercept of a fit should be interpreted in context.
\end{itemize}

\end{document}
